\scanpage{42}\mspage{29}

est dans $C$. Enfin, si $W$\inserted{$\subset\mathrm{Fl}\,M$} est comme dans 2.1.$3^\circ$, et si, $\varnothing_I\longrightarrow X$ est dans $W_{I^{\unsure{o|\circ}}}$ i.e. les $X_i$ sont $\unsure{W|\acute W|\widetilde W}$-coasphériques, alors \deleted{\unsure{$\varnothing$}} $X_\bullet$ est dans $M(I^o)_{TC}$ (où $TC=C\cap W$), et pour tt ouvert $J$ de $I$, $X_{\deleted{I}J}=\underset{i\in\deleted{I}J}{\unsure{\underrightarrow{\lim}|\underline{\lim}}}\,X_i$ est dans $\inserted{M}_{TC}$, \inserted{donc $\unsure{X_j|X_J}$ est $W$-coasphérique,} et pour tt couple d'ouverts $K\subset J$ la flèche de transition $\tau_{K,J}$ est dans $TC$.

Signalons enfin le complément suivant au théorème 2.1., qui résultera de sa démonstration :

\uline{Corollaire} 2.5. Supposons que l'espace topologique associé à $I$ (\inserted{dont} les ouverts sont les ouverts de $I$, et les $I_{/i}$ forment une base d'ouverts) soit \uline{noethérien}, i.e. toute suite croissante d'ouverts est stationnaire. Alors \unsure{toute|toutes} les conclusions dans le th. 2.1. et ses corollaires restent \unsure{valables}, en supposant seulement la stabilité de $C$ ou de $TC=C\cap W$ par composition simple (au lieu de composition transfinie), en plus des deux autres stabilités.

% archive sheet number: 41

